\begin{table}[H]
\centering
\caption{Summary Statistics: Pre-Treatment Period (2015--2019)}
\label{tab:sumstats}
\begin{adjustbox}{max width=\textwidth}
\begin{threeparttable}
\begin{tabular}{l cc cc cc}
\toprule
 & \multicolumn{2}{c}{Mining} & \multicolumn{2}{c}{Non-Mining} & \multicolumn{2}{c}{Difference} \\
\cmidrule(lr){2-3} \cmidrule(lr){4-5} \cmidrule(lr){6-7}
 & Mean & SD & Mean & SD & Diff. & $p$-value \\
\midrule
Total crime & 1327.7 & (5039.1) & 817.3 & (3290.7) & 510.5 & 0.001 \\
Homicide & 21.3 & (85.9) & 10.5 & (50.2) & 10.8 & 0.000 \\
Robbery & 273.1 & (1172.3) & 231.4 & (1169.4) & 41.7 & 0.269 \\
Extortion & 2.8 & (10.1) & 3.2 & (14.2) & -0.4 & 0.258 \\
Domestic violence & 144.4 & (592.4) & 76.6 & (318.9) & 67.8 & 0.000 \\
\midrule
Municipalities & \multicolumn{2}{c}{178} & \multicolumn{2}{c}{2286} & & \\
Municipality $\times$ years & \multicolumn{2}{c}{890} & \multicolumn{2}{c}{11,430} & & \\
\bottomrule
\end{tabular}
\begin{tablenotes}[flushleft]\footnotesize
\item \textit{Notes:} Pre-treatment period is 2015--2019. Mining municipalities are those that received Fondo Minero allocations according to SEDATU 2017 distribution data. Crime counts are annual municipality-level totals from SESNSP. Standard deviations in parentheses. $p$-values from two-sample $t$-tests for equality of means.
\end{tablenotes}
\end{threeparttable}
\end{adjustbox}
\end{table}

